\documentclass[11pt,a4paper]{article}

\usepackage[margin=1in]{geometry}
% \usepackage{microtype}  % disabled — needs scalable fonts; not critical for readability
\usepackage{booktabs}
\usepackage{hyperref}
\usepackage{url}
\usepackage{listings}
\usepackage{xcolor}
\usepackage{amsmath}
\usepackage{amssymb}
\usepackage{graphicx}
\usepackage{caption}

\hypersetup{
    colorlinks=true,
    linkcolor=blue!50!black,
    citecolor=blue!50!black,
    urlcolor=blue!50!black,
}

\lstset{
    basicstyle=\ttfamily\small,
    breaklines=true,
    columns=fullflexible,
    keepspaces=true,
    frame=single,
    framerule=0pt,
    backgroundcolor=\color{black!3},
}

\title{MoE-Infer: A Rust/Metal Engine for Streaming-Expert MoE\\
Inference on Apple Silicon}
\author{Claude (Anthropic)\thanks{Implementation, measurement, and writeup
by Claude (Opus 4.7). Project direction, architecture decisions, and
prior MoE-Infer baseline by Khanh Nguyen.} \and Khanh Nguyen}
\date{May 2026}

\begin{document}

\maketitle

\begin{abstract}
We present MoE-Infer, an inference engine for Mixture-of-Experts language
models that streams expert weights from SSD on demand and runs all
compute through hand-tuned Metal shaders. Starting from the
\texttt{flash-moe} reference~\cite{flashmoe} (a C engine that demonstrated
35B-parameter MoE inference on a laptop), we (i) identified and fixed
two correctness bugs in the baseline (a norm-shift mismatch between
zero-centered and gated RMSNorm, and an MTP embedding indexing bug),
(ii) rewrote the engine in Rust with stricter type safety and a
``deferred command buffer'' design that interleaves adjacent layers into
$N{+}1$ Metal command buffers, and (iii) added a \emph{batched prefill
path} that processes $N$ prompt tokens per layer instead of per token,
collapsing per-token MoE op2 commits into one commit per layer and
loading each unique expert once per layer instead of once per token.
On Qwen3.6-35B-A3B-bq4 we measure $2.7\times$ peak prefill speedup vs
the token-serial path on Apple M4, and sequential-generation throughput
improves from $8$~tok/s (\texttt{flash-moe}) to $10$~tok/s (MoE-Infer)
on M1 Pro for the same model. We also report several systematically
negative results: low-rank expert factorization fails because Qwen3.6
experts are rank-saturated past layer 0, and rotated quantization
(QuIP\#-style) does not help against the existing group-wise BF16-scale
INT4 scheme.
\end{abstract}

\section{Introduction}

Mixture-of-Experts models route each token through only a subset of
``expert'' subnetworks, activating $K \ll E$ of $E$ total experts per
layer. For Qwen3.6-35B-A3B this is $K{=}8$ of $E{=}256$ experts per
layer. Inference cost scales with active parameters (3B), but storage
scales with total parameters (35B). Both flash-moe~\cite{flashmoe} and
this work exploit that gap by keeping only currently-active expert
weights in GPU memory and streaming them from SSD on demand.

Qwen3.6-35B-A3B's architecture~\cite{qwen3} is hybrid: of 40 transformer
blocks, 10 use full multi-head attention with KV cache (every fourth
layer) and 30 use Gated DeltaNet linear attention (fixed-size recurrent
state). All 40 blocks have a MoE FFN with 256 routed experts plus one
shared expert.

This report covers MoE-Infer's three contributions: correctness fixes
relative to the \texttt{flash-moe} baseline (\S\ref{sec:correctness}),
sequential-generation optimizations that yield $\sim 25\%$ throughput
gain on M1 Pro (\S\ref{sec:seq}), and a new batched-prefill engine
that yields up to $2.7\times$ on M4 (\S\ref{sec:batched}). We also
document negative-result experiments
(\S\ref{sec:negative}) and telemetry-driven analysis of where time
goes in the pipeline (\S\ref{sec:telemetry}).

\section{Background}

\paragraph{flash-moe.} The reference C engine~\cite{flashmoe} streams
4-bit expert weights via \texttt{pread} into Metal buffers, dispatches
GPU compute through three sequential command buffers per layer, and
runs attention math on the CPU. It reported 8~tok/s on Qwen3.5-35B-A3B
on Apple M1 Pro. Its key insight: SSD read bandwidth ($\sim 5\!-\!10$~GB/s)
suffices to stream $K \cdot S_{\text{expert}}$ bytes per layer per
token when overlapped with GPU compute.

\paragraph{Model.} Per token, layer:
\begin{itemize}
\item Full-attn or linear-attn block: norms, projections, attention
    compute, residual.
\item MoE block: router gate produces softmax over 256 experts; top-$K$
    selected; per-active-expert (gate, up, down) projections in INT4;
    one shared expert in INT4; weighted combine with residual.
\end{itemize}
Per-layer expert file: $E$ experts $\times$ ($S_{\text{gate}} +
S_{\text{up}} + S_{\text{down}}$) bytes. For Qwen3.6 with
$\texttt{moe\_intermediate}{=}512$, $\texttt{hidden}{=}2048$, INT4 +
BF16 scales/biases: $S_{\text{expert}} \approx 1.77$~MB.

\section{Correctness Fixes}\label{sec:correctness}

Two real bugs in the baseline, both found by integrating against
HuggingFace transformers as the reference (\texttt{verify\_nway.py},
\texttt{helpers/verify\_vs\_original.py}).

\paragraph{Norm-shift overgeneralization.} Qwen3.6 uses
\emph{zero-centered} RMSNorm in most places (the quantization pipeline
adds $+1$ to the norm weight, and the kernel computes $x \cdot (w + 1)$)
but \emph{gated} RMSNorm inside the DeltaNet block ($x \cdot w$, no
shift). The original suffix-match rule
\texttt{is\_norm\_key} ended on \texttt{.norm.weight}, which matched
the gated norm too. The engine then multiplied gated-norm output by
$w{+}1$ instead of $w$. Fixed by an explicit suffix list and an
explicit exclusion of \texttt{.linear\_attn.norm.weight}.

\paragraph{MTP embedding index bug.} The MTP draft path's CPU
\texttt{embed\_token\_cpu} indexed the scales and biases arrays with
the wrong base offset (\texttt{s[g]} instead of \texttt{s\_row[g]}),
producing garbage embeddings for any token ID $> 0$. Fixed by adding
the per-row pointer indirection.

Both fixes were verified against the original BF16 model via cosine
similarity on logits: the engine now matches the dequantized
HuggingFace reference at cosine $> 0.9999$ on the stripped 4-layer
validation model.

\section{Sequential-Generation Engine}\label{sec:seq}

We rewrite \texttt{flash-moe}'s engine in Rust with several design
changes that together yield $\sim 25\%$ throughput improvement on M1
Pro for the same model and hardware.

\paragraph{$N{+}1$ command buffers for $N$ layers.} The forward pass
has exactly one hard sync point per layer: the router decision (the CPU
needs the gate scores before issuing \texttt{pread} for the chosen
experts). \texttt{flash-moe} commits and waits three times per layer.
We collapse this to one wait per layer by encoding ``op2 of layer~$N$
plus op1 of layer~$N{+}1$'' into a shared command buffer, so each
commit covers two layers' worth of GPU work and the sync barrier
amortizes.

\paragraph{Maximal Metal kernel coverage.} The original CPU attention
math (Q/K head-norm, RoPE, KV cache writes, SDPA, sigmoid gate) is
moved to GPU kernels. The engine now has 35 specialized kernels vs
flash-moe's 26.

\paragraph{Block-aware quantization (BQ4).} A mixed-precision scheme:
attention/router weights stay in BF16, \texttt{lm\_head} in INT8
(per-channel symmetric), expert weights INT4 with per-group BF16
scale/bias at group size 64. Costs $\sim 5\%$ throughput vs an
all-INT4 build but recovers most of the quantization gap; see
\S\ref{sec:negative} for the analysis.

\paragraph{Application-level expert cache.} A small shared LRU on
\texttt{(layer, expert\_idx)} avoids re-streaming hot experts. Size is
configurable (default off); 32-entry cache is the sweet spot on M1 Pro
(memory cost $\sim 54$~MB, recovers $\sim 10\%$ throughput). On M4,
GPU compute outruns SSD bandwidth even without the cache, so the cache
is dead weight there.

\paragraph{Measured.} On M1 Pro (10-core CPU, 14-core GPU),
Qwen3.5-35B-A3B-4bit, full 40-layer model, 32-token prompt,
100-token greedy generation:

\begin{center}
\begin{tabular}{lr}
\toprule
Engine & tok/s \\
\midrule
\texttt{flash-moe} (C reference) & 8.0 \\
MoE-Infer all-INT4 & 10.0 \\
MoE-Infer BQ4 & 9.5 \\
\bottomrule
\end{tabular}
\end{center}

\section{Batched-Prefill Engine}\label{sec:batched}

For multi-token prompts (chat prefill, RAG context loading), the
token-serial engine above leaves much of the GPU idle. Each step
dispatches matvecs against weight matrices that are large relative to
the input vector, and the GPU's wide compute units are barely
saturated. \emph{Batched prefill} fixes this by processing $N$ tokens
together at the same point in the layer stack, sharing weight reads
and collapsing per-token command-buffer commits into one per layer.

\subsection{Optimization passes}

We landed seven distinct optimizations on the branch
\texttt{batched-prefill}. Each preserves bit-near-identical logits
($\text{max\_diff} \sim 2 \cdot 10^{-5}$, top-1 match on all sampled
tokens) compared to the token-serial path.

\begin{enumerate}
\item \textbf{Batched matvec kernels} (\texttt{matvec\_bf16\_n},
    \texttt{matvec\_int8\_n}, \texttt{dequant\_matvec\_4bit\_n}):
    Metal kernels taking $[N, \text{in\_dim}]$ input and producing
    $[N, \text{out\_dim}]$ output against a shared weight matrix. The
    grid is linearized
    ($\texttt{tgid\_flat} = \texttt{row\_tile} + n \cdot \texttt{num\_row\_tiles}$)
    because Metal disallows mixing scalar and vector grid attributes
    within a single kernel.
\item \textbf{Causal batched SDPA} (\texttt{attn\_sdpa\_causal\_n}):
    prefill-side analogue of the single-token fused SDPA. One
    threadgroup per ($i$, $h$) where $i \in [0, N)$ is a new token
    and $h$ a query head; online softmax over positions
    $[0, \text{past\_pos} + i]$; causal mask is implicit in the loop
    bound. Paired with \texttt{kv\_cache\_append\_n} which writes $N$
    K/V vectors at positions $[\text{past\_pos}, \text{past\_pos} + N)$
    in a single dispatch.
\item \textbf{\texttt{op1\_full\_batched}}: full-attention layer's
    pre-MoE work for $N$ tokens in one command buffer.
    Per-token-position kernels (Q/K head-norm + RoPE) loop within the
    same encoder with offset-into-batched-buffer addressing; the rest
    use batched matvecs.
\item \textbf{\texttt{op1\_linear\_batched}}: linear-attention layer's
    pre-MoE work for $N$ tokens in one command buffer. DeltaNet's
    recurrence ($\texttt{state}_t = f(\texttt{state}_{t-1}, x_t)$) is
    sequential by definition, but Metal serializes per-token
    dispatches via implicit barriers on the shared
    \texttt{conv\_state} and \texttt{delta\_state} buffers. A tiny
    \texttt{buffer\_copy\_f32} GPU kernel copies between batched and
    single-token buffers within the encoder (CPU memcpy is not
    interleavable with encoding because shared-storage writes are
    only visible at commit).
\item \textbf{Single-commit batched MoE op2}: all $N$ tokens' op2
    dispatches encoded into one command buffer per layer, with
    \texttt{encode\_post\_expert\_at} reading per-token slices of
    \texttt{BatchedFullBuffers} via offset arithmetic. Per-token
    \texttt{combine\_params} live in
    \texttt{bufs.combine\_params\_n[$ti \cdot 10$..]}.
\item \textbf{Unique-expert pool}: the largest single win. Without
    it, $N \cdot K$ expert pre-reads per layer; with it, only the
    \emph{unique} experts in the batch (up to $E$) are pre-read once
    each, then multiple tokens' \texttt{Routing.expert\_data} clones
    point to the same pool slot. For random routing at $N{=}128,
    K{=}8, E{=}256$: $\sim 200\!-\!256$ unique vs 1024 pre-reads.
    Memory cost: $E \cdot S_{\text{expert}} \approx 450$~MB on
    Qwen3.6.
\item \textbf{GEMM-tiled \texttt{matvec\_bf16\_gemm\_n}}: processes
    \texttt{NCOLS\_PER\_TG}${=}4$ tokens per threadgroup with $K$-axis
    tiling (\texttt{TILE\_K}${=}256$). Weight reads are shared
    across the 4 tokens in a threadgroup, reducing weight bandwidth
    $\sim 4\times$ on the BF16 matvecs (attention projections, gate,
    shared-expert projections).
\end{enumerate}

\subsection{Results}

Measured on Apple M4, Qwen3.6-35B-A3B-bq4 full model. Each row is
the minimum of 5 measured runs after one warmup forward.

\begin{center}
\begin{tabular}{rrrrrr}
\toprule
$N$ & seq tok/s & bat tok/s & speedup & seq ms & bat ms \\
\midrule
8   & 6.51  & 11.66 & 1.79$\times$ & 1229  & 686   \\
16  & 6.29  & 12.01 & 1.91$\times$ & 2543  & 1333  \\
32  & 6.77  & 18.45 & \textbf{2.73}$\times$ & 4728  & 1735  \\
64  & 6.29  & 14.92 & 2.37$\times$ & 10170 & 4289  \\
128 & 6.89  & 17.82 & 2.59$\times$ & 18582 & 7182  \\
\bottomrule
\end{tabular}
\end{center}

Notable: speedup grows with $N$ (the unique-expert-pool savings
compound), and absolute throughput climbs from $\sim 7$~tok/s to
$12\!-\!18$~tok/s. Compared to the same engine without batched op2
(\textit{i.e.,} per-token pool, per-token MoE commits), the
unique-expert pool alone accounted for a $+30\!-\!40\%$ jump at large
$N$ (Table~\ref{tab:incremental}).

\begin{table}[h]
\centering
\caption{Incremental contribution of each optimization (speedup vs
    token-serial baseline, $N{=}128$).}
\label{tab:incremental}
\begin{tabular}{lr}
\toprule
Cumulative optimizations & Speedup \\
\midrule
batched lm\_head only                          & 1.05$\times$ \\
+ \texttt{op1\_full\_batched}                  & 1.16$\times$ \\
+ single-commit op2 (per-token pool)           & 1.62$\times$ \\
+ batched op1 for linear-attn                  & 1.72$\times$ \\
+ GEMM-tiled \texttt{matvec\_bf16\_n}          & 1.88$\times$ \\
+ unique-expert pool                           & \textbf{2.59}$\times$ \\
\bottomrule
\end{tabular}
\end{table}

\section{Negative-Result Experiments}\label{sec:negative}

We tried two well-known weight-compression techniques. Both failed on
Qwen3.6 in ways that are themselves informative.

\paragraph{Low-rank expert factorization.} We hypothesized that expert
weight matrices $W \in \mathbb{R}^{m \times n}$ might admit a low-rank
approximation $W \approx U V^\top$ with $r \ll \min(m, n)$, reducing
both storage and pread bandwidth. We measured singular value decay
across all 41 MoE layers (script:
\texttt{helpers/analyze\_expert\_lowrank.py}). Findings:
\begin{itemize}
\item Layer 0 alone is meaningfully low-rank: rank 282 of 512 captures
    99\% variance, $\sim 28\%$ storage saving.
\item Layers 1--40 are rank-saturated: rank $> 485$ of 512 needed for
    99\% variance, and the low-rank decomposition is actually
    \emph{larger} than the dense INT4 matrix it replaces ($+21\%$
    storage).
\item Verdict: low-rank decomposition fails on 40 of 41 layers. The
    experts have been trained to use their full rank.
\end{itemize}

\paragraph{Rotated quantization (QuIP\#-style).} We applied random
orthogonal rotation to expert weight columns and re-quantized at INT2
/ INT3 / INT4, measuring per-matrix output cosine similarity against
the original. Findings: rotation gives $\leq 0.2\%$ cosine lift,
\emph{only} at INT2 where the absolute baseline is already broken
(cosine $\sim 0.91$). At INT3 and INT4 the rotation is a wash or
slight loss. The per-group BF16-scale quantization (group size 64)
already does the local outlier adaptation that rotation would unlock
in coarser schemes. Reverting to INT3 alone drops per-matrix cosine
from 0.996 to 0.982, which compounds catastrophically across $\sim 960$
matrix ops per forward, so INT3 is not viable either.

\section{Telemetry Analysis}\label{sec:telemetry}

We instrumented the engine to attribute time across GPU wait, disk
I/O, and CPU routing (script: \texttt{helpers/bench\_io.py}). On M4,
Qwen3.6 full model, $N{=}63$ tokens, 5 runs:

\begin{center}
\begin{tabular}{lrr}
\toprule
Stage & ms (mean) & share \\
\midrule
Engine total          & 9248 & 100\% \\
\quad GPU op1 wait    & 4654 & 50.3\% \\
\quad Expert routing+I/O & 3577 & 38.7\% \\
\qquad pread (disk only) & 3571 & 38.6\% \\
\quad Routing CPU (softmax+topk) & 8.5 & 0.1\% \\
\bottomrule
\end{tabular}
\end{center}

Two observations: (i)~CPU routing is negligible, validating the
decision to not GPU-side the routing (an earlier experiment that did
GPU routing regressed throughput by adding 40 Metal dispatches per
token). (ii)~GPU and disk are nearly balanced ($\sim 50\%$ each), so
\emph{either} bottleneck dominates wall time, not their sum --- the
$N{+}1$ command-buffer pipelining hides disk I/O behind GPU compute.

The expert cache (\S\ref{sec:seq}) is dead weight on M4 because pread
already fits inside the GPU shadow. On M1 Pro the GPU is slower and
disk is slower; pread becomes the critical path; the cache helps.
Hardware-specific.

\section{Conclusion}

MoE-Infer demonstrates that careful systems work on the GPU/CPU/SSD
boundary yields large practical wins for streaming-expert MoE
inference. The headline numbers are $25\%$ sequential throughput
improvement over the reference engine (10 vs 8 tok/s, M1~Pro) and up
to $2.7\times$ prefill speedup via batching on M4. Both come from the
same theme: minimize the number of CPU/GPU synchronization barriers
and maximize the work covered per barrier. The negative-result
experiments around weight compression (low-rank, rotated quantization)
support a complementary observation: on a well-trained MoE the
weights are not casually compressible, and the engineering leverage
lives in I/O scheduling and dispatch fusion rather than fewer bits per
weight.

The code is at
\url{https://github.com/fbundle/moe_infer}; the batched-prefill engine
is on the \texttt{batched-prefill} branch.

\begin{thebibliography}{9}

\bibitem{flashmoe}
  D. Han. \emph{flash-moe: MoE inference for laptops}.
  \url{https://github.com/danielhanchen/flash-moe}, 2024.

\bibitem{qwen3}
  Qwen Team. \emph{Qwen3 Technical Report}. arXiv:2505.09388, 2025.

\bibitem{quipsharp}
  A. Tseng, J. Chee, Q. Sun, V. Kuleshov, C. De~Sa.
  \emph{QuIP\#: Even Better LLM Quantization with Hadamard Incoherence
  and Lattice Codebooks}. arXiv:2402.04396, 2024.

\bibitem{turboquant}
  A. Zandieh, M. Daliri, M. Hadian, V. Mirrokni.
  \emph{TurboQuant: Online Vector Quantization with Near-optimal
  Distortion Rate}. arXiv:2504.19874, 2025.

\end{thebibliography}

\end{document}
